\subsection{Historical background}
Euler's evaluation of the even zeta values,
\[
 \zeta(2k)=(-1)^{k+1}\frac{B_{2k}(2\pi)^{2k}}{2(2k)!},
\]
places their arithmetic in a different setting from that of the odd values;
see \cite[\S25.6]{DLMF}. Ap\'ery's 1978 discovery, published in
\cite{Apery}, established the irrationality of $\zeta(3)$ through specially
structured rational approximations. Beukers subsequently gave an integral
proof \cite{Beukers}. These arguments illustrate a recurrent principle:
one must construct an expression that is both analytically small and
arithmetically constrained, while also proving that it does not vanish.
The last condition is indispensable, even when strong denominator bounds
and decay estimates are available.

Rivoal \cite{Rivoal00}, and Ball and Rivoal \cite{BR}, established that
infinitely many odd zeta values are irrational. More quantitatively, their
work gives a logarithmically growing lower bound for the dimension over
$\Q$ of the span of $1,\zeta(3),\zeta(5),\ldots,\zeta(s)$ as the odd
integer $s$ tends to infinity. This is a linear-independence statement,
stronger than a corresponding count of irrational entries. Fischler's
Bourbaki account \cite{Fischler04} explains the connections among rational
approximation, hypergeometric series, and the irrationality criteria used
in this development.

Explicit short lists pose a related but distinct problem. Rivoal's
nine-value result \cite{Rivoal02}, announced in 2001, concerns
$\zeta(5),\zeta(7),\ldots,\zeta(21)$. Zudilin proved in 2001 that at least
one of $\zeta(5),\zeta(7),\zeta(9),\zeta(11)$ is irrational
\cite{Zudilin01}. His later systematic treatment \cite{Zudilin04} develops
hypergeometric constructions and their arithmetic, including the removal
of common prime factors. Integral and series representations of such
linear forms are also developed by Nesterenko \cite{Nesterenko03}.

Subsequent work has pursued both more elementary proofs and stronger
asymptotic counts. Zudilin \cite{Zudilin18} obtained an elementary proof
that one of the odd values from $\zeta(5)$ through $\zeta(25)$ is irrational,
using the prime number theorem and Stirling's formula among its analytic
ingredients. Fischler, Sprang and Zudilin \cite{FSZ} proved a lower bound
$2^{(1-\varepsilon)\log s/\log\log s}$ for the number of irrational odd
zeta values up to $s$, for sufficiently large $s$. Lai and Yu \cite{LY}
improved this to $(1.192507\ldots-\varepsilon)\sqrt{s/\log s}$.
Lai's subsequent refinement \cite{Lai26} gives
$1.284\sqrt{s/\log s}$ among $\zeta(3),\zeta(5),\ldots,\zeta(s-1)$ for
sufficiently large even $s$, combining the elimination method with
Zudilin's prime-divisor improvement. These are counts of irrational
values; they should not be confused with lower bounds of the same size
for the dimension of their rational span.

The fixed set studied here is $\{\zeta(7),\zeta(9),\zeta(11),\zeta(13)\}$.
It excludes $\zeta(5)$ and includes $\zeta(13)$, so the result of
\cite{Zudilin01} alone does not imply the desired conclusion. Nor does an
asymptotic counting theorem automatically settle this particular short
interval. The comparison concerns the supports of the assertions, not a
claim to subsume those earlier theorems. The discussion above records relevant precedents rather than an
exhaustive priority survey.

\subsection{The determinant construction and its verification}
The starting point is a factorial-normalized rational function, in the
tradition of \cite{BR,Zudilin04}. Its fourth Taylor coefficient shifts the
surviving odd pole orders $3,5,7,9$ to the arguments $7,9,11,13$.
The present argument uses a growing moment matrix rather than requiring
one prescribed scalar linear form to be nonzero. A finite Pearson
classification and a truncated Hankel extension force its rank to be
within four of its size, uniformly under arbitrary complex substitutions
for the four target zeta values. One then selects an almost full nonzero
minor. Uniform estimates for every minor of that size ensure that the
selection creates no unproved size condition.

On the analytic side, a tilted contour and an explicit bounded-root basis
give a uniform upper exponent. On the arithmetic side, local square-minor
valuations are combined by a finite dual inequality; rational boxes control
all floor values on whole parameter cells. The use of $p$-adic valuations
here concerns rational numbers and the product formula. The theorem is
about the classical real zeta function, not a $p$-adic zeta function.

This version incorporates the completed Round~29 reduction and final audit
\cite{Project}. Lean~4 \cite{Lean4}, Mathlib \cite{Mathlib}, and a pinned
prime-number-theorem development \cite{PNT} check the mathematical
implications from explicit finite certificates to the classical zeta
conclusion. Eight isolated admissions supply the finite numerical families;
separate programs check their data using exact rational arithmetic and Arb
ball arithmetic \cite{Arb}. Thus the result has a completed formal reduction
with externally checked finite inputs. It is not an admission-free Lean
formalization. Section~\ref{sec:verification} and
Appendix~\ref{app:reproduction} specify the boundary and reproduction steps.

\subsection{Statement and rates}
